\subsection{Locked target-trial analysis and final inference}

The primary analysis was restricted to verified HR-positive/HER2-negative patients who were alive and eligible at a day-180 landmark. The treatment contrast compared initiation of hormone therapy during days 0--180 with no initiation by day 180. The target estimand was the treated-minus-control difference in 730-day post-landmark restricted mean survival time (RMST) in the treatment-overlap population.

The locked estimator used five-fold cross-fitting repeated over 20 prespecified nuisance partitions. Propensity scores were estimated using the refitted regularized Stage-30 specification. Censoring was handled with a cross-fitted discrete-time regularized logistic model and a censoring-survival floor of $G_{\min}=0.10$. Arm-specific ridge regressions were fitted to the IPCW-RMST pseudo-outcome, and predicted conditional RMST values were bounded to the admissible interval $[0,730]$ days. Patient-level ATO-AIPW scores were averaged across the 20 partitions.

Before the final resampling analysis, the cohort, estimand, learners, censoring rule, partition seeds, bootstrap seeds, and interval definitions were cryptographically locked. Statistical uncertainty was quantified using 300 ordinary patient-bootstrap repetitions, with all nuisance models refitted within every bootstrap partition. Copies of the same sampled patient were constrained to the same nuisance fold. The prospectively designated primary interval was the 95\% percentile patient-bootstrap interval; basic and studentized intervals were retained as sensitivity analyses.
